% !TeX root = ../tfg.tex
% !TeX encoding = utf8
%
%*******************************************************
% Capítulo 2. Estado del arte
%*******************************************************

\chapter{Estado del arte}\label{cap:estado-arte}

\section{Clasificación de histopatología con aprendizaje profundo}\label{sec:ea-histopatologia}

% Puntos que debería cubrir:
% - Imágenes de lámina completa (WSI), extracción de parches y tinción H&E.
% - Ajuste fino de redes preentrenadas en histopatología (Kandel y Castelli, 2020).
% - Graduación de Gleason y variabilidad entre observadores; conjunto de datos de próstata del
%   grupo (Ayensa-Jiménez et al.) y PANDA (Bulten et al., 2022).
% - Arquitecturas: CNN (ResNet, EfficientNetV2, ConvNeXt) y transformers (ViT, Swin).
\pendiente{estado del arte de la clasificación histopatológica}

\section{Detección de datos fuera de distribución}\label{sec:ea-ood}

\subsection{Taxonomía de la detección OOD generalizada}\label{sec:ea-taxonomia}

% Puntos que debería cubrir: AD, ND, OSR, OOD y detección de valores atípicos (Yang et al., 2024).
\pendiente{taxonomía}

\subsection{Cambio semántico y desplazamiento de covariables}\label{sec:ea-shift}

\pendiente{cambio semántico frente a desplazamiento de covariables}

\subsection{Near-OOD y far-OOD}\label{sec:ea-near-far}

% Puntos que debería cubrir: definición de near-OOD y far-OOD; ubicación de los conjuntos de este
% trabajo (NCT-CRC como near-OOD, DTD como far-OOD).
\pendiente{near-OOD y far-OOD}

\subsection{Criterio conservador en medicina}\label{sec:ea-conservador}

\pendiente{criterio conservador}

\subsection{Métodos excluidos}\label{sec:ea-excluidos}

\pendiente{métodos excluidos}

\section{Explicabilidad}\label{sec:ea-explicabilidad}

\subsection{El efecto Clever Hans}\label{sec:ea-clever-hans}

\pendiente{efecto Clever Hans}

\subsection{Explicaciones locales y conceptuales}\label{sec:ea-locales-conceptuales}

\pendiente{explicaciones locales frente a conceptuales}

\subsection{LRP, CRP y explicaciones glocales}\label{sec:ea-crp}

\pendiente{LRP, CRP y explicaciones glocales}

\endinput
